\appendix
\section{Proof Details}\label{app:B}
\subsection*{Lemma~\ref{lem:1} (Client Drift Bound)}

\noindent\textbf{Step 1 --- Drift recursion.}
\(\omega_{t,e}^k - \omega_t = -\eta_t\sum_{i=0}^{e-1}g_{t,i}^k\).
Cauchy--Schwarz:
\(\|\omega_{t,e}^k - \omega_t\|^2 \le \eta_t^2 e\sum_{i=0}^{e-1}\|g_{t,i}^k\|^2\).
Taking \(\mathbb E_t\) and weighted summation:
\(D_e \le \eta_t^2 e\sum_{i=0}^{e-1}\sum_k p_k\,\mathbb E_t\|g_{t,i}^k\|^2\).

\noindent\textbf{Step 2a --- Noise decomposition.}
\(g_{t,i}^k = \nabla F_k(\omega_{t,i}^k) + \varepsilon_{t,i}^k\)
(\(\mathbb E_t[\varepsilon] = 0\), \(\mathbb E_t\|\varepsilon\|^2 \le \sigma_l^2\)).
By \(\|a+b\|^2 \le 2\|a\|^2 + 2\|b\|^2\):
\(\mathbb E_t\|g_{t,i}^k\|^2 \le 2\sigma_l^2 + 2\|\nabla F_k(\omega_{t,i}^k)\|^2\).

\noindent\textbf{Step 2b --- Pullback.}
\(\nabla F_k(\omega_{t,i}^k) = \nabla F_k(\omega_t) + (\nabla F_k(\omega_{t,i}^k) - \nabla F_k(\omega_t))\).
\(L\)-smoothness:
\(\|\nabla F_k(\omega_{t,i}^k)\|^2 \le 2\|\nabla F_k(\omega_t)\|^2 + 2L^2\|\omega_{t,i}^k - \omega_t\|^2\).

\noindent\textbf{Step 2c --- Combine.}
\(\mathbb E_t\|g_{t,i}^k\|^2 \le 2\sigma_l^2 + 4\|\nabla F_k(\omega_t)\|^2 + 4L^2\,\mathbb E_t\|\omega_{t,i}^k - \omega_t\|^2\).

\noindent\textbf{Step 2d --- Weighted summation.}
\(\sum_k p_k\,\mathbb E_t\|g_{t,i}^k\|^2
 \le 2\sigma_l^2 + 4\sum_k p_k\|\nabla F_k(\omega_t)\|^2 + 4L^2 D_i\).
Using the heterogeneity assumption~\ref{asmp:2}:
\(\le A_t + 4L^2 D_i,\; A_t \triangleq 2\sigma_l^2 + 4\kappa^2 + 4M\|\nabla F(\omega_t)\|^2\).

\noindent\textbf{Steps 3--5 --- Solve via prefix maximum.}
Substituting yields \(D_e \le \eta_t^2 e^2 A_t + 4L^2\eta_t^2 e\sum_{i=0}^{e-1}D_i\).
Defining \(D_{\max,e} \triangleq \max_{0\le j\le e} D_j\):
\(D_{\max,e} \le \eta_t^2 e^2 A_t + 4L^2\eta_t^2 e^2 D_{\max,e}\).
Under \(8L^2\eta_t^2 E^2 \le 1\), we have \(4L^2\eta_t^2 e^2 \le 1/2\),
so \(D_{\max,e} \le 2\eta_t^2 e^2 A_t\). \(\square\)

\subsection*{Lemma~\ref{lem:2} (Gradient Bias and Second Moment)}

\noindent\textbf{Bias term \(\delta_t\):}
\(\delta_t = \frac{1}{E}\sum_{e,k} p_k(\nabla F_k(\omega_{t,e}^k) - \nabla F_k(\omega_t))\).
By Jensen's inequality and \(L\)-smoothness:
\(\mathbb E_t\|\delta_t\|^2 \le \frac{L^2}{E}\sum_{e=0}^{E-1} D_e
 \le L^2 D_{\max,E} \le 2L^2\eta_t^2 E^2 A_t\).\;
This establishes Eq.\,\eqref{eq:30}.

\noindent\textbf{Second moment \(\mathbb E_t\|g_t\|^2\):}
\(\mathbb E_t\|g_t\|^2 \le A_t + 4L^2 D_{\max,E} \le (1+8L^2\eta_t^2 E^2)A_t\).
Substituting \(A_t = 2\sigma_l^2 + 4\kappa^2 + 4M\|\nabla F\|^2\):
\begin{align*}
  (1+8L^2\eta_t^2 E^2)A_t
  &= 2(1+8L^2\eta_t^2 E^2)\sigma_l^2
     + 4(1+8L^2\eta_t^2 E^2)\kappa^2 \notag\\
  &\quad + 4M(1+8L^2\eta_t^2 E^2)\,\|\nabla F\|^2.
\end{align*}
Using \(\eta_t \le \bar\eta\) and the definition of \(c_g\) (Eq.\,\eqref{eq:13}):
\(\frac{1}{2}c_g \ge 2(1+8L^2\bar\eta^2 E^2)\),
\(c_g \ge 4(1+8L^2\bar\eta^2 E^2)\), and
\(c_g \ge 4M(1+8L^2\bar\eta^2 E^2)\).
The three terms are bounded by \(\frac{c_g}{2}\sigma_l^2\), \(c_g\kappa^2\), and \(c_g\|\nabla F\|^2\),
summing to \(\le c_g(\sigma_l^2 + \kappa^2 + \|\nabla F\|^2)\). \(\square\)

\subsection*{Lemma~\ref{lem:3} (Single-Step Descent)}

\noindent\textbf{Step 1 --- Inner product.}
\(\mathbb E_t\langle\nabla F, \Delta_t\rangle
 = -\eta_t E\langle\nabla F, \mathbb E_t[g_t]\rangle
   - \eta_t E\langle\nabla F, m_t\rangle
   + \mathbb E_t\langle\nabla F, r_t\rangle\).
Using \(\mathbb E_t[g_t] = \nabla F + \delta_t\):
\begin{align*}
  \text{(A)}\;& -\eta_t E\langle\nabla F, \nabla F + \delta_t\rangle \notag\\
              &\qquad = -\eta_t E\|\nabla F\|^2 - \eta_t E\langle\nabla F, \delta_t\rangle\\
  \text{(B)}\;& -\eta_t E\langle\nabla F, m_t\rangle,\qquad
  \text{(C)}\; \mathbb E_t\langle\nabla F, r_t\rangle.
\end{align*}
Applying Young's inequality (with \(\alpha = 1/2\) for (A), \(\alpha = 1/2\) for (B), \(\alpha = \eta_t E/4\) for (C)):
\begin{align*}
  \text{(A)} &\le -\eta_t E\|\nabla F\|^2 + \tfrac{\eta_t E}{4}\|\nabla F\|^2 \notag\\
             &\qquad + \eta_t E\|\delta_t\|^2
              = -\tfrac{3\eta_t E}{4}\|\nabla F\|^2 + \eta_t E\|\delta_t\|^2,\\
  \text{(B)} &\le \tfrac{\eta_t E}{4}\|\nabla F\|^2 + \eta_t E\|m_t\|^2,\\
  \text{(C)} &\le \tfrac{\eta_t E}{8}\|\nabla F\|^2 \notag\\
            &\qquad + \tfrac{2}{\eta_t E}\,\mathbb E_t\|r_t\|^2.
\end{align*}
Summing:
\begin{align*}
  \mathbb E_t\langle\nabla F, \Delta_t\rangle
  &\le -\tfrac{3}{8}\eta_t E\|\nabla F\|^2 \quad + \eta_t E\|\delta_t\|^2 \notag\\
  &\quad + \eta_t E\!\cdot\! B_\mu^2
     + \tfrac{2}{\eta_t E}\,\mathbb E_t\|r_t\|^2.
\end{align*}

\noindent\textbf{Step 2 --- Quadratic term.}
\(\|\Delta_t\|^2 \le 3\eta_t^2 E^2\|g_t\|^2 + 3\eta_t^2 E^2\|b_t\|^2 + 3\|r_t\|^2\).
Multiplying by \(L/2\), taking \(\mathbb E_t\), and substituting bounds:
\begin{align*}
  \frac{L}{2}\mathbb E_t\|\Delta_t\|^2
  &\le \frac{3L}{2}\eta_t^2 E^2 c_g(\sigma_l^2 + \kappa^2 + \|\nabla F\|^2) \\
  &\quad + \frac{3L}{2}\eta_t^2 E^2(B_\mu^2 + B_V^2)
     + \frac{3L}{2}\,\mathbb E_t\|r_t\|^2.
\end{align*}

\noindent\textbf{Step 3 --- Merging.}
Substituting \(\eta_t E\|\delta_t\|^2 \le 8L^2\eta_t^3 E^3(\sigma_l^2 + \kappa^2 + M\|\nabla F\|^2)\) (from Eq.\,\eqref{eq:30}):
\begin{align*}
  \mathbb E_t[F(\omega_{t+1})]
  &\le F(\omega_t)
     - \tfrac{3}{8}\eta_t E\|\nabla F\|^2
     + \tfrac{3L}{2}\eta_t^2 E^2 c_g\|\nabla F\|^2 \notag\\
  &\quad + 8L^2\eta_t^3 E^3 M\|\nabla F\|^2 \notag\\
  &\quad + \bigl(\tfrac{3L}{2}c_g\eta_t^2 E^2 + 8L^2\eta_t^3 E^3\bigr)(\sigma_l^2 + \kappa^2) \\
  &\quad + \tfrac{3L}{2}\eta_t^2 E^2 B_V^2
     + \eta_t E B_\mu^2
     + \tfrac{3L}{2}\eta_t^2 E^2 B_\mu^2 \notag\\
  &\quad + \bigl(\tfrac{2}{\eta_t E} + \tfrac{3L}{2}\bigr)\,\mathbb E_t\|r_t\|^2.
\end{align*}

\noindent\textbf{Step 4 --- Coefficient absorption.}
The coefficient of \(\|\nabla F\|^2\) is
\(-\frac{3}{8}\eta_t E + \frac{3L}{2}\eta_t^2 E^2 c_g + 8L^2\eta_t^3 E^3 M
 = -\frac{3}{8}\eta_t E + \eta_t E \cdot \eta_t E L \cdot (\frac{3}{2}c_g + 8L\eta_t E M)\).
Under \(\eta_t E L \le c_0 \le \frac{1}{16(1+c_g)}\) and \(c_g \ge 4M\),
one can verify \(\frac{3}{2}c_g + 8L\eta_t E M \le \frac{7}{2}c_g\),
so the coefficient satisfies
\(\le -\frac{3}{8}\eta_t E + \eta_t E \cdot c_0 \cdot \frac{7}{2}c_g
   \le -\frac{\eta_t E}{8}\)
(since \(\frac{7}{2}c_0 c_g < \frac{1}{4}\)).

\noindent\textbf{Step 5 --- Constants \(c_1, c_2\).}
Using \(8L^2\eta_t^3 E^3 \le 8L c_0 \eta_t^2 E^2\) to reduce order:
\((\frac{3L}{2}c_g + 8L c_0)\eta_t^2 E^2(\sigma_l^2 + \kappa^2)
   + \frac{3L}{2}\eta_t^2 E^2 B_V^2
 \le c_1\eta_t^2 E^2(\sigma_l^2 + \kappa^2 + B_V^2)\),
with \(c_1 = \frac{3L}{2}c_g + 8L c_0 + 3L\).
The \(B_\mu^2\) term:
\(\eta_t E + \frac{3L}{2}\eta_t^2 E^2
 = \eta_t E + \frac{3}{2}(\eta_t E L)\eta_t E
 \le (1 + \frac{3}{2}c_0)\eta_t E\),
yielding \(c_2 = 2 + 3c_0\). \(\square\)

% =================================================================
\subsection{Complete Proofs of Theorems}
% =================================================================

\subsection*{Theorem~\ref{thm:1} (Surrogate Objective Convergence)}

\noindent\textbf{Step 1 (Constant setup).}
From Proposition~\ref{prop:2}: \(\|m_t\|^2 \le B_\mu^2 = 4G^2\),
\(\mathbb E_t\|\zeta_t\|^2 \le B_V^2 = 4G^2\).
From Proposition~\ref{prop:3}: \(\|r_t\|^2 \le D/(t+1)^{2+2\rho}\).

\noindent\textbf{Step 2 (Telescoping sum).}
Take full expectation of Eq.\,\eqref{eq:33} and sum.
Let \(S_T \triangleq \sum_{t=0}^{T-1}\eta_t\).
The five terms accumulate as:
\begin{itemize}[leftmargin=*,nosep]
  \item \(\mathbb E_t[F(\omega_{t+1})] \le \mathbb E_t[F(\omega_t)]
        \;\Longrightarrow\; \mathbb E[F(\omega_T)] \le \mathbb E[F(\omega_0)]\)
        (telescoping cancellation)
  \item \(-\dfrac{\eta_t E}{8}\|\nabla F(\omega_t)\|^2
        \;\Longrightarrow\; -\dfrac{E}{8}\sum_{t=0}^{T-1}\eta_t\,\mathbb E\|\nabla F(\omega_t)\|^2\)
  \item \(c_1\eta_t^2 E^2(\sigma_l^2 + \kappa^2 + B_V^2)
        \;\Longrightarrow\; c_1 E^2(\sigma_l^2 + \kappa^2 + 4G^2)\sum_{t=0}^{T-1}\eta_t^2\)
  \item \(c_2\eta_t E \cdot 4G^2
        \;\Longrightarrow\; 4c_2 G^2 E \cdot S_T\)
  \item \(\bigl(\dfrac{8}{\eta_t E} + \dfrac{3L}{2}\bigr)\mathbb E_t\|r_t\|^2
        \;\Longrightarrow\; \sum_{t=0}^{T-1}\bigl(\dfrac{8}{\eta_t E} + \dfrac{3L}{2}\bigr)\mathbb E\|r_t\|^2\)
\end{itemize}

\noindent\textbf{Step 3 (Rearranging).}
From \(F(\omega_T) \ge F_{\inf}\):
\(\frac{E}{8}\sum\eta_t\,\mathbb E\|\nabla F\|^2
 \le \Delta_F + c_1 E^2(\sigma_l^2 + \kappa^2 + 4G^2)\sum\eta_t^2
    + 4c_2 G^2 E \cdot S_T + \mathcal R_\infty\).

\noindent\textbf{Steps 4--5 (Extract minimum \& series estimates).}
\(\min_{0\le t<T} \le \frac{8}{E S_T}[\Delta_F + c_1 E^2(\cdots)\sum\eta_t^2 + 4c_2 G^2 E \cdot S_T + \mathcal R_\infty]\).
With \(\eta_t = \eta_0/\sqrt{t+1}\):
\(S_T \ge \eta_0\sqrt{T}\;(T\ge 4)\), \(\sum\eta_t^2 \le \eta_0^2(1 + \ln T)\).

\noindent\textbf{Step 6 (Four-term substitution).}

\(\displaystyle\frac{8}{E S_T}\Delta_F \le \frac{8\Delta_F}{E\eta_0\sqrt{T}}\)

\(\displaystyle\frac{8}{E S_T}\,c_1 E^2(\cdots)\sum\eta_t^2
           \le \frac{8c_1 E\eta_0(1+\ln T)}{\sqrt{T}}(\sigma_l^2 + \kappa^2 + 4G^2)\)

\(\displaystyle\frac{8}{E S_T}\cdot 4c_2 G^2 E \cdot S_T = 32 c_2 G^2 E\)
           (\textbf{error floor}, \(S_T\) cancels)
           
\(\displaystyle\frac{8}{E S_T}\mathcal R_\infty \le \frac{8\mathcal R_\infty}{E\eta_0\sqrt{T}}\)


\noindent\textbf{Step 7 (Residual convergence verification).}
\(\mathcal R_\infty
 \le D\Bigl[\frac{8}{E\eta_0}\sum(t+1)^{-3/2-2\rho}
           + \frac{3L}{2}\sum(t+1)^{-2-2\rho}\Bigr] < \infty\)
(since \(\rho > 0\)). \(\square\)

\subsection*{Theorem~\ref{thm:2} (Adaptive Objective Convergence)}

\noindent\textbf{Step 1.}
Proposition~\ref{prop:5} guarantees that Eq.\,\eqref{eq:33} holds for \(\widetilde F_t\).

\noindent\textbf{Step 2.}
From \(\widetilde F_{t+1}(\omega_{t+1}) \le \widetilde F_t(\omega_{t+1}) + \varepsilon_t\) (Assumption~\ref{asmp:4}),
taking full expectation:
\begin{align*}
  \mathbb E[\widetilde F_{t+1}]
  &\le \mathbb E[\widetilde F_t]
     - \tfrac{\eta_t E}{8}\,\mathbb E\|\nabla\widetilde F_t\|^2
     + c_1\eta_t^2 E^2(\cdots) \\
  &\quad + 4c_2 G^2\eta_t E
     + \bigl(\tfrac{8}{\eta_t E} + \tfrac{3L}{2}\bigr)\mathbb E\|r_t\|^2
     + \varepsilon_t.
\end{align*}

\noindent\textbf{Step 3 (Telescoping sum).}
Summing from \(t = 0\) to \(T-1\):
\(\frac{E}{8}\sum\eta_t\,\mathbb E\|\nabla\widetilde F_t\|^2
 \le \widetilde\Delta_0 + \sum\varepsilon_t
    + c_1 E^2(\cdots)\sum\eta_t^2
    + 4c_2 G^2 E \cdot S_T + \mathcal R_\infty\).
Dividing by \(S_T\) and taking the minimum:
\(\frac{8\sum\varepsilon_t}{E S_T} \le \frac{8\mathcal E_\infty}{E\eta_0\sqrt{T}}\);
the error floor \(32 c_2 G^2 E\) remains unchanged. \(\square\)

%%%%%%%%%%%%%%%%%%%%%%%%%%%%%%%%%%%%%%%%%%%%%%%%%%%%%%%%%%%%%%

\bibliographystyle{IEEEtran}
\bibliography{bibfile}

%%%%%%%%%%%%%%%%%%%%%%%%%%%%%%%%%%%%%%%%%%%%%%%%%%%%%%%%%%%%%%

\begin{IEEEbiography}[{\includegraphics[width=1in,height=1.25in,clip,keepaspectratio]{fig/picture/Cuiyongsheng.jpg}}]{Yongsheng Cui}
	(Graduate Student Member, IEEE)
	is currently pursuing the Ph.D. degree in the Department of Electronic Engineering, Tsinghua University, Beijing, China.
	His research interests include deep space communications, satellite constellation design,
	machine learning and satellite TT\&C systems.
\end{IEEEbiography}

\begin{IEEEbiography}[{\includegraphics[width=1in,height=1.25in,clip,keepaspectratio]{fig/picture/Chentaiyi.jpg}}]{Taiyi Chen}
	(Student Member, IEEE)
	received the Master degree in telecommunication engineering from the Department of Electronic Engineering,
	Tsinghua University, Beijing, P.R. China.
	His research interests include Wireless Communication Systems for satellite constellations,
	channel coding, and deep space communications.
\end{IEEEbiography}

\begin{IEEEbiography}[{\includegraphics[width=1in,height=1.25in,clip,keepaspectratio]{fig/picture/Haoran Xie-eps-converted-to.pdf}}]{Haoran Xie}
	(Member, IEEE)
	is a researcher at the Electric Power Research Institute, China Southern Power Grid.
	His research interests include application of optimization, machine learning,
	and statistical theory.
	He has served as the Chair for IEEE IWCMC 2022/2023/2024 Symposium and IEEE VTC 2022-Fall Symposium.
\end{IEEEbiography}

\begin{IEEEbiography}[{\includegraphics[width=1in,height=1.25in,clip,keepaspectratio]{fig/picture/YafengZhan-eps-converted-to.pdf}}]{Yafeng Zhan}
	(Member, IEEE)
	received B.S.E.E. and Ph.D.E.E. degrees from the Department of Electronic Engineering, Tsinghua University, Beijing, China,
	in 1999 and 2004, respectively.
	He is currently a Professor with the Beijing National Research Center for Information Science and Technology,
	Tsinghua University. His current research interests include satellite TT\&C systems,
	communication signal processing and deep space communications.
	He serves as an editor of \textit{Chinese Space Science and Technology} and the \textit{Journal of Deep Space Exploration}.
\end{IEEEbiography}
